\section{Docker Networks}

Docker Networks controlam como os containers se comunicam entre si e com o mundo externo. Ao instalar o Docker, três redes são criadas automaticamente:

\begin{verbatim}
docker network ls
NETWORK ID     NAME      DRIVER    SCOPE
17e324f45964   bridge    bridge    local
6ed54d316334   host      host      local
7092879f2cc8   none      null      local
\end{verbatim}

\subsection{Bridge (padrão)}

O driver \textbf{bridge} é o padrão. Ele cria uma rede virtual privada no host, e os containers conectados a ela podem se comunicar entre si. Containers em redes bridge diferentes não se comunicam diretamente.

\begin{verbatim}
# Containers se conectam ao bridge padrão automaticamente
docker run -d --name app1 nginx
docker run -d --name app2 nginx
\end{verbatim}

Na rede bridge padrão, containers só se comunicam por endereço IP --- \textbf{não por nome}. Para resolver isso, usa-se uma rede bridge customizada (user-defined), que possui resolução DNS automática:

\begin{verbatim}
# Criar rede customizada
docker network create minha-rede

# Containers na rede customizada podem se comunicar por nome
docker run -d --name web --network minha-rede nginx
docker run -d --name api --network minha-rede node-app

# Dentro do container "api", é possível acessar "web" pelo nome:
docker exec api curl http://web:80
\end{verbatim}

Redes bridge customizadas são superiores à bridge padrão porque oferecem resolução DNS automática entre containers, melhor isolamento e configuração mais flexível.

Para expor um container ao mundo externo, usa-se o mapeamento de portas com \texttt{-p}:

\begin{verbatim}
docker run -d -p 8080:80 --network minha-rede nginx
\end{verbatim}

\subsection{Host}

O driver \textbf{host} remove o isolamento de rede entre o container e o host. O container compartilha diretamente a pilha de rede da máquina host --- não há NAT nem mapeamento de portas.

\begin{verbatim}
docker run -d --network host nginx
# Nginx fica acessível na porta 80 do host diretamente
\end{verbatim}

Quando usar host:
\begin{itemize}
    \item Performance máxima de rede (sem overhead de NAT).
    \item Aplicações que precisam escutar em muitas portas dinâmicas.
    \item Ferramentas de monitoramento de rede.
\end{itemize}

Quando \textbf{não} usar host:
\begin{itemize}
    \item Quando isolamento de rede é necessário.
    \item Quando múltiplos containers precisam usar a mesma porta.
    \item Em produção onde isolamento é requisito de segurança.
\end{itemize}

\textbf{Observação}: o modo host só funciona nativamente no Linux. No macOS e Windows, containers rodam dentro de uma VM, então \texttt{--network host} conecta à rede da VM, não à máquina real.

\subsection{None}

O driver \textbf{none} desabilita completamente a rede do container. O container recebe apenas uma interface \textit{loopback} (\texttt{127.0.0.1}) e não consegue se comunicar com nada externo.

\begin{verbatim}
docker run --network none alpine ping 8.8.8.8
# ping: sendto: Network unreachable
\end{verbatim}

É útil para containers que processam dados localmente e não precisam de nenhum acesso à rede, aumentando a segurança e o isolamento.

\subsection{Comandos úteis para redes}

\begin{verbatim}
docker network ls                          # listar redes
docker network create <nome>               # criar rede
docker network inspect <nome>              # detalhes da rede
docker network connect <rede> <container>  # conectar container a rede
docker network disconnect <rede> <cont.>   # desconectar
docker network rm <nome>                   # remover rede
docker network prune                       # remover redes não utilizadas
\end{verbatim}

\subsection{Resumo}

\begin{center}
\begin{tabular}{|l|l|l|}
\hline
\textbf{Driver} & \textbf{Isolamento} & \textbf{Comunicação} \\
\hline
Bridge & Rede virtual isolada & Containers na mesma rede; porta para host \\
Host & Sem isolamento & Compartilha rede do host diretamente \\
None & Total & Sem rede (apenas loopback) \\
\hline
\end{tabular}
\end{center}
